%% SYNTHETIC TEST UNIT generated by writing/tools/make_outline_book.py from FINAL_OUTLINE_v2 -- not book text
\chapter{有效场论方法：SCET}\label{ch:scet}
本单元为测试用的合成文本（由 FINAL\_OUTLINE\_v2 自动生成），陈述内容从略。

\section{第 1 节（测试）}\label{sec:scet-modes}

\begin{definition}[模式]\label{def:scet-modes}
（测试用合成陈述：内容从略。）
\uses{def:ir-scalings}
\end{definition}

\begin{proposition}[场的标度 \tagR]\label{prop:scet-scalings}
（$R_{\mathrm{def}}$：由算符定义直接推出。）
（测试用合成陈述：内容从略。）
\uses{def:scet-modes, def:prelim-feynman-rules}
\end{proposition}

\begin{proposition}[领头阶共线夸克拉氏量 \tagR]\label{prop:scet-lagrangian}
（$R_{\mathrm{PT}}$：在重整化微扰论中对红外正规化的部分子态严格。）
（测试用合成陈述：内容从略。）
\uses{prop:scet-scalings, def:prelim-qcd}
\end{proposition}

\section{第 2 节（测试）}\label{sec:scet-lagrangian}

\begin{proposition}[共线 Wilson 线 \tagR]\label{prop:scet-wilson}
（$R_{\mathrm{PT}}$：在重整化微扰论中对红外正规化的部分子态严格。）
（测试用合成陈述：内容从略。）
\uses{prop:scet-lagrangian, def:prelim-wilson}
\end{proposition}

\begin{example}[$W_n$ 的树图匹配]\label{ex:scet-wn-tree}
（测试用合成陈述：内容从略。）
\uses{prop:scet-wilson}
\end{example}

\begin{proposition}[共线与超软规范对称性 \tagR]\label{prop:scet-gauge}
（$R_{\mathrm{def}}$：由算符定义直接推出。）
（测试用合成陈述：内容从略。）
\uses{def:scet-modes, prop:qcd-covariance, lem:prelim-brst-nilpotent}
\end{proposition}

\section{第 3 节（测试）}\label{sec:scet-decoupling}

\begin{theorem}[超软退耦（BPS 场重定义） \tagR]\label{thm:scet-decoupling}
（$R_{\mathrm{def}}$：由算符定义直接推出。）
（测试用合成陈述：内容从略。）
\uses{prop:scet-gauge, prop:prelim-wilson-covariance}
\end{theorem}

\begin{example}[BPS 场重定义的树图演示]\label{ex:scet-decoupling-tree}
（测试用合成陈述：内容从略。）
\uses{thm:scet-decoupling}
\end{example}

\begin{remark}[$Y_n$ 的方向与 $i0$]\label{rem:scet-y-direction}
（测试用合成陈述：内容从略。）
% uses: none
\end{remark}

\section{第 4 节（测试）}\label{sec:scet-scet2}

\begin{definition}[零箱减除]\label{def:scet-zero-bin}
（测试用合成陈述：内容从略。）
\uses{def:method-subtraction, lem:prelim-scaleless}
\end{definition}

\begin{proposition}[SCET$_{\rm II}$ 与软 Wilson 线 \tagR]\label{prop:scet-scet2}
（$R_{\mathrm{PT}}$：在重整化微扰论中对红外正规化的部分子态严格。）
（测试用合成陈述：内容从略。）
\uses{def:scet-modes, prop:scet-wilson}
\end{proposition}

\begin{proposition}[快度发散的处理与等价性 \tagA]\label{prop:scet-rapidity-rg}
（测试用合成陈述：内容从略。）
\uses{prop:scet-scet2, def:tmd-collins, prop:tmd-evolution}
\end{proposition}
\begin{proofsketch}测试。\end{proofsketch}
\paragraph{地位与缺口}测试：见文献。

\section{第 5 节（测试）}\label{sec:scet-glauber}

\begin{proposition}[Glauber 算符 \tagA]\label{prop:scet-glauber-operators}
（测试用合成陈述：内容从略。）
\uses{def:scet-modes, prop:scet-scet2}
\end{proposition}
\begin{proofsketch}测试。\end{proofsketch}
\paragraph{地位与缺口}测试：见文献。

\begin{remark}[Glauber SCET 与光锥序证明的对照]\label{rem:scet-glauber-dictionary}
（测试用合成陈述：内容从略。）
% uses: none
\end{remark}

\begin{proposition}[硬匹配系数 \tagR]\label{prop:scet-hard-matching}
（$R_{\mathrm{PT}}$：在重整化微扰论中对红外正规化的部分子态严格。）
（测试用合成陈述：内容从略。）
\uses{prop:sudakov-oneloop}
\end{proposition}

\section{第 6 节（测试）}\label{sec:scet-thrust}

\begin{theorem}[推力分布的因子化 \tagA]\label{thm:scet-thrust}
（测试用合成陈述：内容从略。）
非微扰地位：\cref{asm:intro-bridge}。
\uses{thm:scet-decoupling, prop:scet-hard-matching, def:scet-zero-bin, prop:method-glauber-deformation, thm:ir-leading-regions, asm:intro-bridge}
\end{theorem}
\begin{proofsketch}测试。\end{proofsketch}
\paragraph{地位与缺口}测试：本陈述的地位由\cref{prop:method-glauber-deformation}、\cref{lem:coll-final-poles}、\cref{asm:intro-bridge}、\cref{asm:appN-envelope}决定。

\begin{proposition}[单圈因子与 RG 一致性 \tagR]\label{prop:scet-thrust-oneloop}
（$R_{\mathrm{PT}}$：在重整化微扰论中对红外正规化的部分子态严格。）
（测试用合成陈述：内容从略。）
\uses{prop:scet-hard-matching, def:prelim-plus}
\end{proposition}

\begin{example}[单圈 $H$、$J$、$S_T$ 合成 QCD 奇异项]\label{ex:scet-thrust-oneloop}
（测试用合成陈述：内容从略。）
\uses{prop:scet-thrust-oneloop}
\end{example}

\section{第 7 节（测试）}\label{sec:scet-rg}

\begin{proposition}[领头对数指数 \tagA]\label{prop:scet-ll}
（测试用合成陈述：内容从略。）
\uses{thm:scet-thrust, prop:scet-thrust-oneloop, prop:prelim-beta}
\end{proposition}
\begin{proofsketch}测试。\end{proofsketch}
\paragraph{地位与缺口}测试：本陈述的地位由\cref{thm:scet-thrust}决定。

\begin{remark}[SCET 证明的前提]\label{rem:scet-assumptions}
（测试用合成陈述：内容从略。）
% uses: none
\end{remark}
